\documentclass[11pt, a4paper]{article}
\input{preamble.tex}
\switchlang{en}
\begin{document}

\title{1981 -- Xerox Star Mechanical Mouse}
\date{}
\maketitle
\selectlanguage{english}

In April 1981, Xerox Corporation announced the Xerox Star 8010 Information System, considered the first office computer, as well as the first window-based graphical shell with icons, folders, and two-button mouse \cite{wiki}. The mouse was available in two variants: optical mouse, invented by Xerox specially for those computers, and mechanical mouse featuring an older ball-based Xerox Alto mouse design from Jack Hawley, co-author of the original 1973 Xerox patent for mouse with two inclined wheels \cite{pat}. Both  mice shared the same topcase  (fig. \ref{fig:XeroxStarMechanicalPic}) - a further evolutionary development of the Xerox Alto Mouse visual style (the first mass-produced mouse in the history of computers). Most probably, mechanical mouse was mostly used during the initial period of Xerox Star presence when the optical mice were not yet produced in enough quantities: Howley's design itself wasn't too reliable according to the quotes given in \cite{mouses}: they picked up dirt, quickly became gummed, and no longer controlled the cursor, being actually the main reason to invent the optical alternative.

\begin{figure}[h]
    \centering
    \includegraphics[scale=0.7]{1981_xerox_star_mechanical_mouse/pic_30.jpg}
    \caption{Xerox Star Mechanical Mouse}
    \label{fig:XeroxStarMechanicalPic}
\end{figure}

\section*{Exterior}

The Xerox Star mouse body had become slightly smaller in comparison with the original Alto mouse, preserving the original rectangular shape in more elegant manner. The upper part is made of beige plastic, and is still derived from a parallelepiped, widening toward the bottom. Its front and back sides are convex and noticeably inclined, and vertical edges are rounded, which substantially reduces the original association with a ``box''. On the upper side of the case are two large, oblong, rounded buttons made of the darker plastic that mostly follow the shape of Alto mouse buttons (fig. \ref{XeroxStarMechanicalTopAndBottom}).

\begin{figure}[h]
    \centering
    \includegraphics[scale=0.5]{1981_xerox_star_mechanical_mouse/top_30.jpg}
    \includegraphics[scale=0.5]{1981_xerox_star_mechanical_mouse/bottom_30.jpg}
    \caption{Xerox Star Mechanical Mouse, top and bottom views}
    \label{fig:XeroxStarMechanicalTopAndBottom}
\end{figure}

Bottom side is entirely made of metal (figure \ref{fig:XeroxStarMechanicalTopAndBottom}). Rotation is registered by a smooth steel ball in the center, while two smaller balls act as feet to minimize friction. A removable ring that allows you to remove the ball to remove collected debris is not yet provided in this model, so complete disassembly is necessary for cleaning.

\section*{Size}

The mouse could have been smaller and lower in size if it had not used the standard Alto mouse body. As it is, it has the typical dimensions for mechanical mice of the first half of the 80s (figure \ref{fig:XeroxStarMechanicalSize}).

\begin{figure}[h]
    \centering
    \includegraphics[scale=0.4]{1981_xerox_star_mechanical_mouse/size_15.jpg}
    \caption{Xerox Star Mechanical Mouse on a graduated pad with a grid step of 1~cm}
    \label{fig:XeroxStarMechanicalSize}
\end{figure}

\section*{Ergonomics}

In terms of ergonomics and appearance, the Alto Mouse is minimalist. The user experience certainly suffers from the severe ``rectangularity'' of the case. Although this is partly compensated for by the convex oblong buttons, located within the reach of the fingers, the case doesn't provide much palm support (figure \ref{fig:XeroxStarMechanicalHand}).

\begin{figure}[h]
    \centering
    \includegraphics[scale=0.4]{1981_xerox_star_mechanical_mouse/hand_30.jpg}
    \caption{Xerox Star Mechanical Mouse with a human hand model}
    \label{fig:XeroxStarMechanicalHand}
\end{figure}

\section*{Connectivity}

The mouse has connector typical for Xerox Alto (DA-15 in this version), and the same quadrature interface as other Xerox mice.

\section*{Internal design}

Mouse internals are shown in figure \ref{fig:XeroxStarMechanicalInside}. The removable solid protection of the ball is worthy of mention: it requires additional disassembly operations to remove garbage. This mouse uses mechanical encoders  based on a metal contact drum (instead of the disk  more common in subsequent models) and four contacts of different length instead of two, which increase the encoder resolution.

\begin{figure}[h]
    \centering
    \includegraphics[scale=0.8]{1981_xerox_star_mechanical_mouse/inside_30.jpg}
    \caption{Xerox Star Mechanical Mouse disassembled}
    \label{fig:XeroxStarMechanicalInside}
\end{figure}

\begin{thebibliography}{9}
\bibitem {star} Xerox Star: Wikipedia \url{https://en.wikipedia.org/wiki/Xerox_Star}
\bibitem{pat} Transducer for a display-oriented pointing device \url{https://patents.google.com/patent/US3892963A/en}
\bibitem{mouses} Xerox Mice. oldmouse.com \url{https://web.archive.org/web/20210418000634/http://oldmouse.com/mouse/xerox/alto.shtml}

\bibitem {wiki} Xerox Alto: Wikipedia \url{https://en.wikipedia.org/wiki/Xerox_Alto}
\bibitem {vlsi81} R.\,F. Lyon. The Optical Mouse, and an Architectural Methodology for
Smart Digital Sensors // VLSI DESIGN, August 1981. - p. 20--30. \url{https://www.dicklyon.com/tech/OMouse/OpticalMouse-Lyon.pdf}
\bibitem {vlsi82} R.\,F. Lyon, M.\,P. Haeberli. Designing and Testing The Optical Mouse // VLSI DESIGN, January/February, 1982. - p. 20--30. \url{https://www.dicklyon.com/tech/OMouse/DesigningTestingOMouse.pdf}
\bibitem{pad} R.\,F. Lyon The Optical Mouse: Early Biomimetic Embedded Vision / Advances in Embedded Computer Vision, Nov 2014, pp.3-22 \url{https://static.googleusercontent.com/media/research.google.com/ru//pubs/archive/43260.pdf}
\end{thebibliography}
\end{document}
